\section{Introduction}
\label{sec:intro}

Evaluating a model means running it on many benchmarks, and every benchmark costs compute, time and people. Public leaderboards have meanwhile accumulated large model$\times$benchmark score matrices, and the scores of different benchmarks are far from independent. This suggests a concrete budgeted problem. Given a history of $M$ models scored on $N$ benchmarks (with missing cells), a new model can only be run on $k\ll N$ benchmarks: \emph{which $k$, and how accurately can the remaining $N-k$ scores be inferred? How large must $k$ be?}

Existing efficiency work mostly reduces the number of \emph{items inside a benchmark} \citep{vivek2024anchor,polo2024tiny,kipnis2024metabench,perlitz2024efficient}, and work on relations between benchmarks mostly asks \emph{whether} they agree or lie on a low-dimensional manifold \citep{perlitz2024benchbench,ruan2024observational,polo2024sloth}. We study the missing middle: the \emph{choice of whole benchmarks} to run for a new model, with evaluation designed so that the answer is unbiased and comes with uncertainty. We make four contributions.
\begin{enumerate}[leftmargin=1.5em,itemsep=1pt,topsep=1pt]
\item \textbf{A method} (\S\ref{sec:method}): a Gaussian conditional predictor fitted to pairwise-complete covariances of the history, with a greedy variance-reduction rule for choosing $S$ that costs $O(N^2k)$, together with six alternative predictors (mean, averaging, factor model, SVD imputation, nearest neighbours, a blend) and seven alternative selection rules (including random).
\item \textbf{Theory} (\S\ref{sec:theory}, Appendix~\ref{app:proofs}): a closed-form error with an irreducible per-benchmark floor (adaptive selection gains nothing under the Gaussian model); a sandwich $\sum_i g_i/(1+kg_i)\le\text{excess}\le r/(k-r-1)$ for random selection, hence $k=\Theta(r/\epsilon)$; a spectral lower bound valid for \emph{any} selection; and the $k/(M-k-1)$ price of unstructured regression on a finite history. All are verified numerically.
\item \textbf{Evidence on seven real matrices} (three nested subsets of MTEB, Open LLM Leaderboard, two HELM Lite versions, LiveBench), under random-model, new-organisation and temporal splits, with bootstrap CIs, and ablations of history size, missingness, benchmark pool and score transform.
\item \textbf{A protocol point}: scoring each selector on the benchmarks \emph{left over} after its own selection inflates the apparent advantage of informative selection; we score all methods on identical held-out targets.
\end{enumerate}

\paragraph{Findings in brief.} On MTEB-main, one well-chosen task gives $R^2=0.52$ (random: 0.39), five give 0.71, ten 0.76, and accuracy then saturates near 0.80 for \emph{all} selectors. Selection matters mostly for $k\le10$ (MSE $-21\%$ at $k=1$, nil at $k=40$), the most commonly evaluated benchmarks are \emph{worse} than random ones, the benefit shrinks with sparse or small histories, in homogeneous pools and on the smallest LLM matrices, and the fitted Gaussian predicts its own held-out error to within $\sim$10\%, which lets the budget be planned from the history.
